% !TEX program = xelatex
% FedCSL paper in NeurIPS format. Requires neurips_2023.sty from:
% https://neurips.cc/Conferences/2023/AuthorInstructions
% 本文含中文，必须使用 XeLaTeX（Overleaf 中请在 Menu -> Compiler 选择 XeLaTeX）。
\documentclass{article}

% 左侧行号：默认投稿模式由 neurips_2023 自动加 lineno。若需关闭行号，改为 \usepackage[preprint]{neurips_2023}；终稿用 [final]
% [nonatbib]：不让模板自带 author-year 版 natbib，改由下面显式载入 numeric 版
\usepackage[nonatbib]{neurips_2023}
% NeurIPS 发表稿（如 FedPCL NeurIPS 2022）使用的数字方括号引用：[1], [2, 3, 4], [5--7]
\usepackage[numbers,sort&compress]{natbib}

% 使用 XeLaTeX/tectonic 编译，去掉 pdflatex 专用的 inputenc / fontenc / CJKutf8
\usepackage{hyperref}
\usepackage{url}
\usepackage{booktabs}
\usepackage{amsfonts}
\usepackage{amsmath}
\usepackage{amssymb}
\usepackage{amsthm}
\usepackage{bm}
\usepackage{nicefrac}
\usepackage{microtype}
\usepackage{xcolor}
\usepackage{graphicx}
\usepackage{multirow}
\usepackage{array}
\usepackage{algorithm}
\usepackage{algpseudocode}
\usepackage{enumitem}
\usepackage{float}

% 中文支持：XeLaTeX + xeCJK + fandol（tectonic 可自动拉取 fandol-fonts）
\usepackage{fontspec}
\usepackage{xeCJK}
% 用 .otf 文件名而不是字体族名，让 tectonic 能从 bundle 按文件解析
\setCJKmainfont[
  BoldFont      = FandolSong-Bold.otf,
  ItalicFont    = FandolKai-Regular.otf,
  BoldItalicFont= FandolKai-Regular.otf
]{FandolSong-Regular.otf}
\setCJKsansfont[BoldFont=FandolHei-Bold.otf]{FandolHei-Regular.otf}
\setCJKmonofont{FandolFang-Regular.otf}

\hypersetup{
    colorlinks=true,
    linkcolor=black,
    citecolor=black,
    urlcolor=blue,
}

\newtheorem{definition}{定义}
\newtheorem{proposition}{命题}
\newtheorem{theorem}{定理}
\newtheorem{remark}{说明}

\title{FedCSL: 面向非独立同分布多变量时间序列的联邦 Shapelet 学习与 OMP 客户端选择}

\author{%
  作者姓名 \\
  单位名称 \\
  \texttt{email@example.com}
}

\begin{document}

\maketitle

\renewcommand{\abstractname}{摘要}

\begin{abstract}
多变量时间序列广泛存在于人体活动识别、工业监测、医疗诊断与设备运维等场景中。这类数据通常天然分散在不同终端、机构或边缘节点，既难以集中汇聚，又普遍存在标签稀缺、分布异构和隐私敏感等问题。传统集中式时间序列表征学习方法在拥有完整数据时能够取得较好性能，但在联邦环境下往往面临客户端漂移严重、跨客户端表示不一致以及训练效率受限等挑战。现有联邦时间序列方法多侧重监督分类、参数平均或系统调度，而对可解释的无监督 shapelet 表征学习与客户端选择的联合研究仍然不足。

针对上述问题，本文提出一种面向非独立同分布多变量时间序列的联邦 shapelet 学习客户端选择框架 FedCSL。该方法以 TimeCSL/CSL 的多尺度 shapelet 表征学习思想为基础，在联邦训练中引入三类关键机制。首先，在客户端选择层面，本文基于正交匹配追踪方法构造客户端重要性估计，通过分析候选客户端局部模型对全局模型的稀疏重构关系，动态确定参与训练的客户端子集，并与 uniform、Oort 和 FedCS 等基线策略进行比较。其次，在表示学习层面，本文在本地实例级对比学习目标之上，引入局部--全局联合对比学习与联合知识蒸馏机制，使客户端模型在保留本地适应性的同时，与全局语义空间保持一致。再次，在多尺度建模层面，本文对不同 shapelet 尺度分别施加对比与蒸馏约束，并结合 STFT 与 ACF 构造周期感知的尺度权重，提高模型对多样时序模式的适配能力。

本文进一步从优化与表示学习角度分析了所提方法的合理性。我们指出，纯粹依赖参数平均的联邦学习在非 IID 条件下难以维持稳定的表示空间，而局部--全局联合蒸馏能够在特征层面抑制客户端漂移；多尺度约束与周期感知加权则能够降低尺度失配带来的冗余优化；基于 OMP 的客户端选择有助于在训练预算有限的情况下优先保留更具代表性的客户端更新。实验部分基于 UEA 多变量时间序列分类数据集以及若干扩展数据集进行评估，并从总体性能、消融实验、非 IID 强度、收敛趋势和可解释性分析等多个方面验证方法有效性。结果表明，FedCSL 在下游分类性能、表示稳定性和训练效率方面均优于多种联邦客户端选择基线与表示学习基线。
\end{abstract}

\noindent\textbf{关键词：} 联邦学习；多变量时间序列；Shapelet；对比学习；客户端选择；正交匹配追踪

\section{引言}

时间序列分析是人工智能和数据挖掘中的重要研究方向，在智能制造、脑电诊断、人体活动识别、金融预测以及工业告警等场景中均具有广泛应用价值。与静态表格数据或单张图像不同，时间序列不仅包含多维观测值，还包含显著的时间依赖、周期结构和跨尺度模式。因此，如何学习具有判别性、可迁移性和可解释性的时间序列表征，一直是该领域的核心问题之一。

近年来，无监督或自监督时间序列表征学习得到快速发展。相比强依赖标签的监督方法，这类方法能够在标签稀缺甚至无标签的条件下，从原始时间序列中挖掘稳定的语义结构。尤其是基于对比学习的方法，通过构造正负样本关系，使模型在特征空间中捕捉时序信号的内在一致性，从而显著提升下游分类、聚类与异常检测性能\cite{ts2vec,tstcc,tfc}；在此基础上，后续工作进一步从掩码建模\cite{simmtm}、时间--频率一致性\cite{tfc}以及软对比关系\cite{softclt}等不同角度扩展了自监督信号，并在通用时序表示层面取得持续提升\cite{timesurl}。TimeCSL/CSL 工作进一步表明，shapelet 作为时间序列分析中具有天然解释性的子序列模式，不仅可以用于任务特定建模，还可以与对比学习结合，形成通用的无监督 shapelet 表征框架\cite{timescsl}。这为后续研究提供了一个非常有价值的方向，即在保持模型可解释性的前提下提升表示学习能力\cite{interpgn,unishape}。

然而，在现实应用中，大量时间序列数据并不适合集中式处理。一方面，医疗、生理监测、设备传感和边缘终端产生的数据通常分散存储于不同客户端；另一方面，这些数据往往具有较强隐私属性，直接上传原始数据会带来隐私泄露和合规风险。联邦学习为此提供了一种自然的解决方案，即只在本地保留原始数据，通过参数或表示级别的信息交换实现多客户端协同训练\cite{fedavg}。为缓解联邦非 IID 场景下的客户端漂移，现有工作普遍在本地目标中加入约束项（如 FedProx 的近端项\cite{fedprox}）或引入模型级对比监督（如 MOON\cite{moon}）。尽管联邦学习在图像、语言和推荐系统中已被广泛研究，但针对多变量时间序列的联邦表示学习，尤其是可解释的 shapelet 表征学习，仍处于相对早期阶段\cite{fedst,fedtsc}。

把集中式的 TimeCSL/CSL 直接迁移到联邦环境并不是一件简单的事情。首先，联邦场景下不同客户端数据往往呈现严重的非独立同分布特性，不同用户、设备和场景对应的数据统计特征、类别分布和时间模式存在明显差异。若仍采用纯粹的局部训练加参数平均，客户端模型很容易过拟合其本地模式，导致全局模型在表示空间中出现漂移和不稳定现象。其次，时间序列本身具有复杂的多尺度结构，不同数据集甚至同一数据集中的不同样本，其主导模式可能出现在不同的时间尺度上。如果在联邦环境下仍对各尺度一视同仁，则会降低 shapelet 学习的针对性。再次，联邦训练通常无法在每轮都使用全部客户端，因此如何从候选客户端中选择更具代表性、更有助于逼近全局目标的参与者，也会显著影响模型训练效率与最终性能。

已有研究从不同角度探讨了这些问题。在联邦时序分类与蒸馏方向，EFDLS 关注多任务时间序列分类中的特征级师生蒸馏与伙伴匹配，强调在不同任务间安全共享相似知识\cite{efdls}；FedX 在跨客户端引入双向知识蒸馏与对比学习，用于无监督联邦表示学习\cite{fedx}；FedU2 则从表示坍缩与客户端表示空间不一致两个角度，给出统一正则与聚合方案\cite{fedu2}。在客户端选择方向，Oort 从统计效用与系统速度的联合视角进行 guided participant selection\cite{oort}；FedCS 类方法则从系统异构和统计异构共同作用的角度出发，优化客户端采样概率\cite{fedcsadaptive}；近期工作进一步将客户端选择问题形式化为稀疏重构任务，提出基于压缩感知的 OMP 采样方案\cite{secompfl}。然而，现有方法大多将“表示学习”“客户端选择”“联邦优化”分开处理：要么偏重系统效率，要么偏重监督分类精度，较少将可解释的 shapelet 表征学习与客户端选择统一到同一框架内。

基于此，本文提出一种面向非 IID 多变量时间序列的联邦 shapelet 学习与 OMP 客户端选择框架 FedCSL。本文的核心思想是：以 TimeCSL/CSL 的多尺度 shapelet 对比学习框架为骨干，在联邦环境下同时处理“谁参与训练”和“如何进行跨客户端表示对齐”两个关键问题。具体而言，本文首先基于正交匹配追踪（Orthogonal Matching Pursuit, OMP）构造客户端选择机制。对于每一轮训练，服务器不再仅依赖固定或均匀采样，而是根据候选客户端局部模型对当前全局模型的稀疏重构能力，动态估计客户端的重要性，并据此更新下一轮客户端采样概率。随后，在被选中的客户端上进行本地多尺度 shapelet 表征学习。为了缓解非 IID 场景下的客户端漂移，本文在本地实例对比学习之外，引入局部--全局联合对比学习（joint contrastive learning）与联合知识蒸馏（joint knowledge distillation），并进一步在每个尺度子空间上执行 scale-wise 的对齐与蒸馏。为了让模型更好地适应不同数据集的频谱与周期结构，本文还基于 STFT 和 ACF 设计周期感知加权机制，对不同 shapelet 尺度进行动态调节。

综上，本文试图回答以下几个问题：第一，如何将可解释的 shapelet 表征学习扩展到联邦非 IID 环境；第二，如何利用 OMP 构造适合联邦时序场景的客户端选择策略；第三，如何通过局部--全局多尺度知识对齐提升全局表示空间的一致性；第四，如何从理论与实验两个层面解释周期感知加权和质量感知聚合的有效性。

本文主要贡献概括如下：
\begin{enumerate}[leftmargin=2.5em]
    \item 提出一种面向联邦多变量时间序列的 shapelet 表征学习与 OMP 客户端选择一体化框架，在保留 TimeCSL/CSL 可解释性的基础上，实现联邦环境下的协同表征学习。
    \item 设计局部--全局 joint CL/KD 与 multi-scale CL/KD 机制，通过特征层面的软对齐抑制非 IID 场景下的客户端漂移，并增强不同尺度子空间的一致性。
    \item 引入基于 STFT 与 ACF 的周期感知多尺度权重机制，使 shapelet 学习能够根据不同数据的主导周期结构进行自适应调整。
    \item 在客户端选择层面，以 OMP 为本文主方法，并与 uniform、Oort、FedCS 等策略进行比较，说明客户端选择与表示学习主干协同设计的重要性。
    \item 从优化与表示学习视角给出方法合理性的理论解释，并在多类时间序列数据集上通过总体性能、消融实验、收敛分析和可解释性实验验证所提方法的有效性。
\end{enumerate}

\section{相关工作}

\subsection{时间序列无监督表示学习}

时间序列无监督表示学习旨在在缺乏标注数据的前提下学习通用时序表示。早期方法通常依赖重构任务、序列预测或自编码器结构，但这类方法往往更关注局部重建误差，而不一定能够直接形成适合下游任务的判别特征\cite{simmtm}。随后，对比学习被广泛引入时间序列分析领域，其基本思想是在不同增强视图之间最大化正样本一致性，同时抑制负样本相似性，从而学习稳定而有区分度的表示。典型工作如 TS2Vec 在分层上下文视图上进行对比\cite{ts2vec}，TS-TCC 分别构造时间对比与上下文对比模块\cite{tstcc}，TF-C 将时间与频率视图的一致性作为自监督信号\cite{tfc}；更进一步，SoftCLT 以软分配缓解了“邻近样本被强制推开”的问题\cite{softclt}，TimesURL 则系统性地重审了正负对构造与 SSCL 损失设计\cite{timesurl}，而 SimMTM 通过掩码时序建模与流形邻居重构扩展了自监督范式\cite{simmtm}。

现有大量方法以 CNN、Transformer 或其变体为主干进行时序建模\cite{fcnbaseline,timesnet}，但这些架构最初并非专为时间序列设计，在处理不同长度、多尺度以及强周期特征时常常存在解释性不足的问题。一条互补的路线是以 shapelet 等可解释子序列模式为核心构建时序分类器\cite{interpgn,unishape}。TimeCSL/CSL 的意义在于将 shapelet 这一对时间序列更自然、更可解释的模式表示与对比学习相结合，使模型能够从多个尺度和多种相似性度量出发提取通用特征\cite{timescsl}。相比纯黑盒表示，shapelet-based representation 具有更强的可解释性，更便于对模型输出进行可视化分析和任务迁移\cite{interpgn}。

\subsection{联邦时间序列学习与知识蒸馏}

联邦时间序列学习研究的核心问题在于：在不共享原始数据的条件下，如何使多个客户端协同训练得到兼顾隐私与性能的模型\cite{fedavg}。现有工作多聚焦于监督分类场景，其中参数平均、个性化联邦学习、蒸馏和迁移学习是几种常见技术路线\cite{fedkdsurvey}。在时序领域，FedST 提出面向 shapelet 变换的安全联邦协议\cite{fedst}，FedTSC 进一步将可解释 TSC 模型嵌入联邦系统\cite{fedtsc}；EFDLS 则将联邦蒸馏引入多任务时间序列分类，通过特征级教师--学生蒸馏和权重匹配策略促进不同任务间知识共享\cite{efdls}。这类工作说明，在时序数据异构较强的环境中，仅靠参数平均并不足以支持有效的跨客户端知识传递，特征层面的蒸馏和语义对齐具有重要意义\cite{moon,fedx}。

跨领域的联邦无监督表示学习也在快速发展。FedCA 最早针对 FURL 提出跨客户端表示对齐以缓解表示不一致\cite{fedca}；FedMoCo 把 MoCo 迁移到医疗图像联邦场景，辅以元数据迁移与自适应聚合\cite{fedmoco}；FedX 通过双向知识蒸馏与局部--全局对比学习学习无偏表示\cite{fedx}；FedU2 指出 FUSL 同时存在表示坍缩和表示空间不一致，给出统一聚合方案\cite{fedu2}。在有监督侧，MOON 引入模型级对比以纠正本地训练\cite{moon}，FedProc 与 FedPCL 以类原型构造对比信号\cite{fedproc,fedpcl}，Reducing non-IID (CDL) 则用对比散度损失抑制跨车辆异构\cite{fedcdl}，并有工作从互信息视角统一讨论联邦对比的合理性\cite{fedmi}。更贴近本文局部--全局蒸馏思路的是 FLESD\cite{flesd}：它利用集成相似度矩阵在服务器侧进行相似度蒸馏，从而在强约束下实现高效联邦表示学习；近期工作 FLPD\cite{flpd} 与 FedRCIL\cite{fedrcil} 则分别从原型相似度蒸馏与增量对比学习出发，给出通信友好的联邦蒸馏方案。此外，FedCOG\cite{fedcog}、MrTF\cite{mrtf}、FedIOCD\cite{fediocd} 与 FedNoRo\cite{fednoro} 等工作从“生成对齐”“转导推理”“输入输出协同蒸馏”和“噪声鲁棒聚合”等角度进一步拓展了联邦蒸馏的边界。

然而，上述方法大多默认任务目标为监督分类或视觉无监督场景，很少直接研究无监督 shapelet 表征学习问题。相比之下，本文更加关注联邦环境下的通用时序表示，即先学习可解释的表示空间，再在该空间上完成下游任务建模。这一思路既继承了 TimeCSL 的优势\cite{timescsl}，也更适合标签稀缺场景\cite{fedsmall}。

\subsection{联邦客户端选择与参与者采样}

在真实联邦环境中，由于带宽、设备异构性和大规模客户端数量的限制，服务器通常无法在每一轮让所有客户端同时参与训练，因此客户端选择成为影响训练效率与模型性能的重要因素。Oort 从系统效率与统计效用的平衡出发，利用客户端最近一次训练所反映的效用信息和设备速度信息进行 guided participant selection，以提高 time-to-accuracy 表现\cite{oort}。另一类方法如 FedCS 及其相关自适应采样研究，则从系统异构与统计异构共同作用的角度，优化客户端采样概率，以缩短达到目标损失所需的总物理时间\cite{fedcsadaptive}。FOLB 通过评估设备更新对全局收敛的预期增益进行加权采样\cite{folb}；Update-Aware Device Scheduling 则进一步在无线边缘环境中，把本地更新的“重要性”与信道条件共同纳入调度\cite{updateaware}；IGPE 将客户端选择问题刻画为子模优化，以在异构系统中兼顾代表性与效率\cite{igpe}；近期研究也开始关注公平性维度，尝试在效率与公平之间进行权衡\cite{fairerfl}。与本文最为相关的是 SecOMPFL\cite{secompfl}，它首次将联邦客户端选择显式建模为稀疏优化问题，并基于 OMP 识别对全局模型贡献冗余或正交的客户端；本文在此思路基础上，把 OMP 系数作为“贡献度”来更新下一轮采样分布，并额外引入 EMA 平滑与最小采样概率以兼顾稳定性与公平覆盖。

上述方法为联邦训练中的参与者选择提供了重要启发，但它们主要从系统调度或梯度优化角度出发，较少考虑所选择客户端对高层语义表示结构的影响。本文认为，在联邦时序表征学习场景下，客户端选择不应仅服务于训练加速，还应服务于全局表示空间的稳定构建。因此，本文基于 OMP 设计客户端选择策略，并将其与 shapelet 表征学习主干协同优化。

\section{问题定义与总体框架}

\subsection{联邦时序场景定义}

设联邦系统由一个中心服务器和 $K$ 个客户端构成。第 $k$ 个客户端持有本地多变量时间序列数据集
\begin{equation}
D_k = \{x_{k,i}\}_{i=1}^{n_k},
\end{equation}
其中 $x_{k,i}\in \mathbb{R}^{C\times T}$，$C$ 表示通道数，$T$ 表示时间长度，$n_k$ 表示客户端 $k$ 的样本规模。总样本规模为
\begin{equation}
N = \sum_{k=1}^{K} n_k.
\end{equation}

在本文中，不同客户端之间的数据分布通常满足非独立同分布条件，即存在
\begin{equation}
P_k(x)\neq P_j(x),\quad k\neq j.
\end{equation}
这种差异不仅体现在类别分布上，也体现在时间尺度、周期模式、幅值范围以及局部噪声结构上。

\subsection{目标定义}

本文的目标不是直接训练一个端到端分类器，而是学习一个联邦共享的 shapelet 编码器 $f(\cdot;w)$，将原始时间序列映射到可解释的表示空间：
\begin{equation}
z = f(x;w)\in \mathbb{R}^{d}.
\end{equation}
得到表示后，再通过下游分类器（如 SVM）完成分类任务评估。该设计与 TimeCSL 的统一分析范式保持一致，即将表示学习与任务求解解耦。

\subsection{总体框架}

FedCSL 的总体流程如图~\ref{fig:framework} 所示，可概括为以下步骤：
\begin{enumerate}[leftmargin=2.5em]
    \item 服务器初始化全局 shapelet 编码器及客户端采样概率分布；
    \item 在第 $t$ 轮开始时，服务器根据当前客户端选择策略，从候选客户端中采样一个参与子集，其中 OMP 为本文方法，uniform、Oort、FedCS 为对照基线；
    \item 被选中客户端接收全局模型，在本地进行多尺度 shapelet 对比学习，并引入局部--全局联合蒸馏与尺度级蒸馏；
    \item 各客户端上传本地模型更新；
    \item 服务器对客户端更新进行聚合，得到新的全局模型；
    \item 周期性感知模块根据 STFT 与 ACF 统计动态调整多尺度损失权重；
    \item 重复上述过程直到训练收敛，并在学习到的全局表示上执行下游评估。
\end{enumerate}

\begin{figure}[H] % 放在正文引用位置附近
    \centering
    \includegraphics[width=0.98\linewidth]{imageall.png}
    \caption{FedCSL 整体框架图。上半部分为服务器端：基于 OMP 的稀疏重构客户端选择（$\min\|y-Ax\|_2^2$）从客户端池中挑出参与子集 $S_t$，经 Quality-Aware Aggregator 聚合得到全局 Shapelet 编码器 $w_g$，下游以 SVM 做表示评测。下半部分为被选中客户端 $k$ 的本地训练：多变量时间序列经数据增强形成两种视图 $x_q,x_k$，送入 8 个尺度（$0.1T$--$0.8T$）的多尺度 Shapelet 编码器；特征对齐节点聚合本地与全局特征，驱动局部对比损失 $\mathcal{L}_{local}$、局部--全局联合对比 $\mathcal{L}_{joint\text{-}CL}$ 与联合蒸馏 $\mathcal{L}_{joint\text{-}KD}$；右下 Period-Aware Local Weighting Module 由 STFT（短--中尺度）与 ACF（中--长尺度）两分支产生尺度权重 $\bm{\pi}$ 反馈到多尺度编码器。}
    \label{fig:framework}
\end{figure}


\section{方法}

\subsection{多尺度 Shapelet 编码器}

本文采用与 TimeCSL/CSL 一致的多尺度 shapelet 表示思想。对于长度为 $T$ 的输入时间序列，在 $\left[0.1T,0.8T\right]$ 范围内均匀选取 $R=8$ 个尺度，每个尺度学习固定数量的 shapelets。不同于只依赖单一相似性度量的编码方式，本文实现同时支持欧氏距离、余弦相似度、互相关以及它们的混合形式。记第 $r$ 个尺度下的 shapelet 集合为 $\mathcal{S}^{(r)}=\{s_m^{(r)}\}_{m=1}^{M_r}$，则一个样本的最终表示可写为
\begin{equation}
z = \left[z^{(1)},z^{(2)},\dots,z^{(R)}\right],
\end{equation}
其中 $z^{(r)}$ 表示输入序列与尺度 $r$ 下所有 shapelets 的最优匹配响应向量。

这种表示方式有两个优势。其一，不同尺度的 shapelets 分别捕捉局部瞬时模式、中程结构和长期依赖，使模型能覆盖更广的时序模式空间。其二，shapelet 作为具体的子序列模式，天然具备解释性，便于后续通过可视化分析理解模型决策依据。

\subsection{本地实例级对比学习}

对每个被选中的客户端，本文首先在本地样本上构造两种增强视图。增强操作包括加噪、裁剪、池化、量化和时间扭曲等时序特有变换，这一构造延续了 TS2Vec/TS-TCC 等时序对比方法的做法\cite{ts2vec,tstcc}。给定输入样本 $x$，记两种增强视图为 $x_q$ 和 $x_k$，对应表示分别为
\begin{equation}
q = f(x_q;w_k),\qquad k = f(x_k;w_k).
\end{equation}
对表示进行归一化后，构造实例级对比学习目标：
\begin{equation}
\mathcal{L}_{\text{local}} = \mathrm{CE}\left(\frac{qk^\top}{\tau}, y\right),
\end{equation}
其中 $\tau$ 为温度系数，$y$ 表示正样本匹配标签。该目标鼓励同一样本的不同增强视图在表示空间中接近，同时与其他样本保持区分。

\subsection{局部--全局联合对比与知识蒸馏}

在联邦非 IID 条件下，仅依赖本地实例级对比学习容易导致各客户端在不同语义坐标系中优化，进而削弱全局模型聚合效果\cite{fedu2}。这一问题在视觉联邦无监督工作中已被广泛讨论：MOON 通过模型级对比把本地表示拉回全局邻域\cite{moon}，FedX 引入跨客户端双向蒸馏学习无偏表示\cite{fedx}，FLESD 则在服务器侧利用集成相似度矩阵完成相似度蒸馏\cite{flesd}。受此启发，本文将上一轮全局模型作为语义教师，记其在两种增强视图上的表示为
\begin{equation}
q_g = f(x_q; w_g),\qquad k_g = f(x_k; w_g).
\end{equation}

在此基础上，引入联合对比项
\begin{equation}
\mathcal{L}_{\text{joint-CL}} = \mathrm{CE}\left(\frac{q_gk^\top}{\tau}, y\right),
\end{equation}
以及联合蒸馏项
\begin{equation}
\mathcal{L}_{\text{joint-KD}} = D_{\mathrm{KL}}(q \,\|\, q_g) + D_{\mathrm{KL}}(k \,\|\, k_g),
\end{equation}
其中 $D_{\mathrm{KL}}$ 表示基于相似度矩阵构造的 KL 散度。通过这两项，本地模型既保留了对本地时序模式的适应性，又被限制在全局共享语义空间附近，从而减小由数据异构引起的客户端漂移。

\subsection{多尺度结构约束与尺度级蒸馏}

对于表示向量 $z=[z^{(1)},\dots,z^{(R)}]$，整体级别的蒸馏并不能保证每个尺度子空间都学得充分。为此，本文将联合对齐细化到每个尺度：
\begin{equation}
\mathcal{L}_{\text{scale-CL}} = \sum_{r=1}^{R}\alpha_r \,
\mathrm{CE}\left(\frac{q_g^{(r)}(k^{(r)})^\top}{\tau}, y\right),
\label{eq:scale-cl}
\end{equation}
\begin{equation}
\mathcal{L}_{\text{scale-KD}} = \sum_{r=1}^{R}\alpha_r \left[
D_{\mathrm{KL}}\left(q^{(r)} \,\|\, q_g^{(r)}\right)+
D_{\mathrm{KL}}\left(k^{(r)} \,\|\, k_g^{(r)}\right)\right].
\end{equation}
其中 $\alpha_r$ 为第 $r$ 个 shapelet 尺度上的损失权重。启用周期感知（\texttt{UseACF}）时，$\alpha_r$ 由第~\ref{sec:period-score} 节在线计算；否则可取均匀权重。

此外，本文保留了来自 CSL/TimeCSL 的结构性约束。设某个尺度下的表示矩阵为 $Z^{(r)}$，其协方差矩阵为
\begin{equation}
C^{(r)} = \frac{\left(Z^{(r)}\right)^\top Z^{(r)}}{B-1},
\end{equation}
则通过约束协方差非对角项，可以抑制不同 shapelet 通道之间的冗余，提升各尺度表示的独立性与判别性。结合多尺度累积统计，本文形成了结构对齐损失 $\mathcal{L}_{\text{structure}}$，以避免表示坍缩和尺度冗余。

\subsection{周期感知加权机制}
\label{sec:period-score}

多变量时间序列在不同数据集、不同通道上的\textbf{主导周期/主频}往往不同；而 shapelet 编码器在 $R$ 个离散长度（实现中取 $R{=}8$，约为 $0.1L,\ldots,0.8L$ 的参考尺度）上并行提取模式。若对所有尺度等权优化，容易在\textbf{与当前数据频谱/周期不匹配}的尺度上浪费梯度。本文在实现中于每个训练步、基于\textbf{当前 batch} 的原始输入 $\mathcal{X}\in\mathbb{R}^{N\times K\times L}$（$N$ 为 batch 大小，$K$ 为通道数，$L$ 为序列长度）调用 \texttt{period\_score}（见 \texttt{utils.py}），构造与 $R$ 个参考尺度对齐的非负权重，再写入式~\eqref{eq:scale-cl} 及尺度级蒸馏中的 $\alpha_r$（代码中另乘常数以调节梯度幅度，不影响相对比例）。

\paragraph{参考尺度与离散锚点。}
取归一化锚点
\begin{equation}
\bm{\rho}=(\rho_1,\ldots,\rho_8)=(0.1,0.2,\ldots,0.8),
\end{equation}
对应以采样点为单位的\textbf{目标尺度}（实现中对每个样本独立计算）
\begin{equation}
\tau_r=\bigl\lfloor \rho_r L \bigr\rfloor,\qquad r=1,\ldots,8.
\end{equation}
这些锚点与多尺度 shapelet 的长度配置一一对应。

\paragraph{基于 STFT 的频域得分（短—中尺度分支）。}
对 batch 内每个样本的至多 $k=\min(9,K)$ 条通道分别计算短时傅里叶变换。实现采用 \texttt{scipy.signal.stft}，采样率归一化为 $f_s{=}1$，窗长 $W=\max(4,\lfloor L/2\rfloor)$，FFT 长度 $N_{\mathrm{fft}}=8W$，帧移为 $h=\lfloor W/4\rfloor$，相邻帧重叠 $W{-}h$。记通道 $c$ 的功率谱为 $P_c(f,t)=|Z_{xx}(f,t)|^2$。为突出\textbf{能量集中频带}，先取功率的 $95\%$ 分位数作为阈值，构造高能量掩膜 $\mathbb{I}[P_c\ge Q_{0.95}(P_c)]$，再定义\textbf{掩膜加权平均频率}
\begin{equation}
\bar{f}_c=\frac{\sum_{f,t} f\cdot P_c(f,t)\cdot \mathbb{I}[\cdot]}{\sum_{f,t} P_c(f,t)\cdot \mathbb{I}[\cdot]+\varepsilon}.
\end{equation}
在 $f_s{=}1$ 的约定下，$\bar{f}_c$ 可视为该通道\textbf{主频}的代理；实现中将其映射为与“周期长度”同量纲的标量 $s_c^{\mathrm{ST}}=1/(\bar{f}_c+\varepsilon)$（$\bar{f}_c$ 过小则置 $0$）。随后用高斯核将标量 $s_c^{\mathrm{ST}}$ 软对齐到各锚点（实现中对 $r=1,\ldots,8$ 均计算，再对通道在 $r$ 上归一化）：
\begin{equation}
u_{c,r}^{\mathrm{ST}}=\exp\!\left(-\frac{\bigl(s_c^{\mathrm{ST}}-\tau_r\bigr)^2}{2\sigma^2}\right),\quad
\sigma=0.1\,L,\quad r=1,\ldots,8,
\end{equation}
对 batch 与通道累加后\textbf{仅保留前四维} $\mathbf{S}^{\mathrm{ST}}\in\mathbb{R}^4$（对应 $\tau_1,\ldots,\tau_4$，较短—中等 shapelet 尺度），与代码 \texttt{total\_scores\_stft[:4]} 一致。图~\ref{fig:period-stft} 预留为\textbf{单通道 STFT 功率谱/时频图}示意，可标注高能量区域与 $\bar{f}_c$ 的几何含义。


\begin{figure}[htbp]
    \centering
    \includegraphics[width=0.9\linewidth]{../../figs/stft_v5_paper_branch.png}
    \caption{周期感知中 STFT 分支的频谱示意（HAR 真实样本：Walking \#2759 / total\_acc\_x）：左图为单通道 STFT 功率谱，白色等高线为 $95\%$ 高能量掩膜，黄色虚线为加权主频 $\bar f_c$；右图将标量 $1/\bar f_c$ 通过高斯核软对齐到 8 个 shapelet 尺度锚点 $\tau_r$，其中前 4 个锚点由 STFT 分支写入 $\mathbf{S}^{\mathrm{ST}}$，后 4 个由 ACF 分支接管。}
    \label{fig:period-stft}
\end{figure}

\paragraph{基于 ACF 的时域滞后得分（中长—长尺度分支）。}
对同一批通道，使用 \texttt{statsmodels} 的样本自相关函数（\texttt{fft=True}，最大滞后 $L{-}1$）得到 $R_c(\ell)$，$\ell=0,\ldots,L{-}1$。实现仅在\textbf{后半段滞后} $\ell\ge \lfloor L/2\rfloor$ 上，用 \texttt{scipy.find\_peaks} 检测峰值，且要求 $R_c(\ell)>\theta_{\mathrm{acf}}$（代码中 $\theta_{\mathrm{acf}}=0.2$）；方差过小的常数序列跳过。每个峰值滞后 $\ell$ 对\textbf{全部八个}锚点产生反距离投票
\begin{equation}
v_{c,r}(\ell)=\frac{1}{|\tau_r-\ell|+\varepsilon},\qquad
\mathbf{v}_c=\sum_{\ell\in\mathcal{P}_c}\bigl(v_{c,1}(\ell),\ldots,v_{c,8}(\ell)\bigr),
\end{equation}
其中 $\mathcal{P}_c$ 为通道 $c$ 上满足条件的峰值集合；再对 $\mathbf{v}_c$ 按通道归一化并在 batch 与通道上累加，得到 $\mathbf{S}^{\mathrm{ACF}}\in\mathbb{R}^8$。实现中\textbf{仅取后四维} $\mathbf{S}^{\mathrm{ACF}}_{5:8}$ 与较长 shapelet 尺度（$\tau_5,\ldots,\tau_8$）对齐，以强调\textbf{长周期/长记忆}结构。图~\ref{fig:period-acf} 给出\textbf{ACF 曲线随滞后 $\ell$} 的示意，标出阈值 $\theta_{\mathrm{acf}}$ 与后半段峰值。

\begin{figure}[htbp]
    \centering
    \includegraphics[width=0.9\linewidth]{../../figs/acf_v5_paper_branch.png}
    \caption{周期感知中 ACF 分支的滞后—投票示意：左图为样本自相关函数 $R_c(\ell)$ 随滞后 $\ell$ 的变化，红色虚线为阈值 $\theta_{\mathrm{acf}}$，后半段（$\ell\ge\lfloor L/2\rfloor$）检出的峰值以圆点标记；右图将这些峰值按反距离投票软对齐到 8 个 shapelet 尺度锚点 $\tau_r$，后 4 个锚点由 ACF 分支写入 $\mathbf{S}^{\mathrm{ACF}}_{5:8}$，与 STFT 分支互补。}
    \label{fig:period-acf}
\end{figure}

\paragraph{双分支融合与损失中的使用。}
将 STFT 分支向量归一化得 $\hat{\mathbf{S}}^{\mathrm{ST}}=\mathbf{S}^{\mathrm{ST}}/\|\mathbf{S}^{\mathrm{ST}}\|_1$，ACF 后四维归一化得 $\hat{\mathbf{S}}^{\mathrm{ACF}}_{5:8}$。融合系数 $\beta\in[0,1]$（实现中对应配置项 \texttt{beta}，调用 \texttt{period\_score(..., alpha=self.beta)}）控制频域与滞后域信任比例：
\begin{equation}
\bm{\pi}=\mathrm{Norm}\!\left(
\beta\cdot \mathrm{concat}\!\bigl(\hat{\mathbf{S}}^{\mathrm{ST}},\mathbf{0}_{4}\bigr)
\;+\;
(1-\beta)\cdot \mathrm{concat}\!\bigl(\mathbf{0}_{4},\hat{\mathbf{S}}^{\mathrm{ACF}}_{5:8}\bigr)
\right),
\label{eq:period-fuse}
\end{equation}
其中 $\mathrm{Norm}(\cdot)$ 表示对 8 维向量做非负归一化使和为 $1$；若各分支全零则退化为均匀权重。实现返回前将 $\bm{\pi}$ 乘以常数 $10$ 以便与损失量级匹配；在 \texttt{train.py} 中，第 $r$ 个尺度的对比/蒸馏项再乘以与 $\pi_r$ 成正比的系数（代码为 $\pi_r\times 5$ 写入 \texttt{precisions}）。因此，\textbf{同一前向 batch} 的频谱与 ACF 统计直接决定该步多尺度损失的相对 emphasis，无需额外标签。

\subsection{OMP 客户端选择}

本文将客户端选择视为一个“从候选客户端更新中稀疏重构全局模型”的问题。这一建模思路与最近将联邦客户端选择显式写为稀疏优化的 SecOMPFL\cite{secompfl} 一致，但本文并不追求与服务器模型的精确重构，而是将 OMP 返回的稀疏系数用作\textbf{客户端重要性代理}，并以此动态更新下一轮采样概率。设当前轮共有 $M$ 个候选客户端，其本地模型参数展平后组成字典矩阵
\begin{equation}
A = [a_1,a_2,\dots,a_M]\in \mathbb{R}^{d_w\times M},
\end{equation}
其中 $a_m$ 表示第 $m$ 个客户端局部模型参数向量，$d_w$ 为参数总维度；记当前全局模型参数展平后为
\begin{equation}
y\in \mathbb{R}^{d_w}.
\end{equation}
则 OMP 客户端选择可以抽象为如下稀疏重构问题：
\begin{equation}
\min_{x}\ \|y-Ax\|_2^2\quad
\text{s.t.}\quad \|x\|_0 \leq s,
\end{equation}
其中 $x\in \mathbb{R}^{M}$ 为稀疏系数向量，$s$ 为允许的非零系数个数。系数绝对值越大，说明对应客户端对重构当前全局模型越重要。

在实现上，本文首先使用 OMP 得到稀疏系数向量 $x$（对负系数取绝对值参与贡献度）。记掩码向量 $\bm{m}\in\{0,1\}^M$ 表示第 $t$ 轮各客户端是否被选中参与训练。对被选中的客户端，定义稳健贡献度
\begin{equation}
\tilde{x}_i = \max\bigl(|x_i|,\varepsilon\bigr),\qquad m_i=1,
\end{equation}
其中 $\varepsilon>0$ 为极小常数，避免 $|x_i|$ 过小者在归一化分配中占比为零；未参与本轮训练的客户端在构造目标分布时仅保留 $p_i^{(t)}$，见下式。

\textbf{目标分布。}在仅根据本轮 OMP 结果构造的“瞬时”目标概率 $\hat{\bm{p}}^{(t)}$ 上，对被选中客户端按贡献度比例重分配其在上一轮所占据的总概率质量；未被选中的客户端目标概率暂保持上一轮取值：
\begin{equation}
\hat{p}_i^{(t)} =
\begin{cases}
\dfrac{\tilde{x}_i}{\sum_{j:m_j=1}\tilde{x}_j}\cdot \displaystyle\sum_{j:m_j=1}p_j^{(t)}, & m_i=1,\\[1.2ex]
p_i^{(t)}, & m_i=0.
\end{cases}
\end{equation}

\textbf{指数移动平均（可调平滑）。}直接将 $\hat{\bm{p}}^{(t)}$ 用作下一轮采样分布会使概率轨迹对单轮 OMP 波动过于敏感。实现中引入\textbf{可调平滑系数} $\eta\in(0,1]$（实现参数 \texttt{ema\_alpha}），对目标分布与上一轮分布做凸组合：
\begin{equation}
\tilde{p}_i^{(t)}=\eta\,\hat{p}_i^{(t)}+(1-\eta)\,p_i^{(t)}.
\end{equation}
$\eta$ 越大，采样分布越紧跟当前 OMP 给出的相对重要性；$\eta$ 越小，分布变化越平缓，有利于训练稳定。$\eta$ 可在实验中通过配置文件或命令行调节（典型范围如 $0.1$--$0.4$）。

\textbf{最小概率截断。}为防止部分客户端长期被边缘化，在归一化前对每个分量施加\textbf{可调下限} $p_{\min}\in(0,1/M)$（实现参数 \texttt{min\_selection\_prob}，默认约 $0.01$）：
\begin{equation}
p_i^{(t+1)}=\frac{\max\bigl(\tilde{p}_i^{(t)},\,p_{\min}\bigr)}{\sum_{j=1}^{M}\max\bigl(\tilde{p}_j^{(t)},\,p_{\min}\bigr)}.
\end{equation}
最终 $\sum_i p_i^{(t+1)}=1$，并据此对下一轮客户端子集进行无放回抽样。相较于纯随机策略，上述设计在强调 OMP 重构代表性的同时，通过 $(\eta,\,p_{\min})$ 在“响应速度—稳定性—公平覆盖”之间提供显式调节手段。

\subsection{联邦聚合}

在客户端选择之后，服务器对被选中客户端上传的模型进行聚合。记第 $k$ 个客户端上传模型为 $w_k$，其样本量为 $n_k$，客户端质量分数为 $q_k$，则可定义聚合权重
\begin{equation}
\omega_k = \frac{n_k q_k}{\sum_{j\in \mathcal{S}_t} n_j q_j},
\end{equation}
其中 $\mathcal{S}_t$ 为第 $t$ 轮被选中的客户端集合。全局更新为
\begin{equation}
w_g^{(t+1)} = \sum_{k\in \mathcal{S}_t} \omega_k w_k^{(t+1)}.
\end{equation}
当不启用质量感知分数时，$q_k$ 可退化为常数，从而回到标准的样本量加权聚合形式。

\subsection{总体优化目标}

综上，本文在客户端 $k$ 上的总体目标可写为
\begin{equation}
\begin{aligned}
\mathcal{L}_k ={}& \gamma\,\mathcal{L}_{\text{local}}
+ \lambda_5\,\mathcal{L}_{\text{structure}} \\
&+ \mathbb{I}_{\text{FedCSL}}\Bigl[
\lambda_1\,\mathcal{L}_{\text{joint-CL}}
+ \lambda_2\,\mathcal{L}_{\text{joint-KD}}
+ \lambda_3\,\mathcal{L}_{\text{scale-CL}}
+ \lambda_4\,\mathcal{L}_{\text{scale-KD}}
\Bigr],
\end{aligned}
\end{equation}
其中 $\mathbb{I}_{\text{FedCSL}}$ 在实现中为 \texttt{algo=fedcsl} 时取 $1$，否则取 $0$，其后各项分别由配置项 \texttt{UseJointCL}、\texttt{UseJointKD}、\texttt{UseScaleCL}、\texttt{UseScaleKD} 控制是否计入。$\mathcal{L}_{\text{structure}}$ 在代码里对应 $\texttt{l3}\times(\mathcal{L}_{\text{cca}}+\texttt{l4}\times\mathcal{L}_{\text{sdl}})$，故式中 $\lambda_5$、$\lambda_{\text{sdl}}$ 与超参 \texttt{l3}、\texttt{l4} 相当。$\lambda_1,\ldots,\lambda_4$ \textbf{不是}互不相关的自由超参：\texttt{train.py} 的 \texttt{update\_CL} 用 $\gamma$、$\zeta=1-\gamma$ 及固定因子 $\frac12$、$\frac{1}{10}$ 对联合项与多尺度项做缩放，且多尺度上 \texttt{UseACF} 时 $\alpha_r$ 与 \texttt{precisions}（含因子 $5$）相乘；联合蒸馏中 \texttt{direct\_kl\_loss} 的两项对 $\zeta$ 的耦合方式亦与书上理想的独立 $\lambda_2$ 略有差别。
若采用 FedProx 形式的参数约束\cite{fedprox}，则还可额外加入
\begin{equation}
\mathcal{L}_{\text{prox}} = \frac{\mu}{2}\|w_k-w_g\|_2^2.
\end{equation}
其中各项损失是否开启，由实验配置文件决定。

\begin{algorithm}[H]
\caption{FedCSL 训练流程}
\label{alg:fedcsl}
\begin{algorithmic}[1]
\State 初始化全局模型 $w_g^{(0)}$、客户端采样概率 $\bm{p}^{(0)}$
\For{$t=0,1,\dots,T-1$}
    \State 根据当前策略选择客户端集合 $\mathcal{S}_t$
    \If{使用 OMP}
        \State 根据候选客户端局部模型与全局模型构造稀疏重构问题
        \State 得到系数向量 $x$，由 $\bm{m}^{(t)}$ 与 $|x|$ 构造目标分布 $\hat{\bm{p}}^{(t)}$
        \State $\tilde{\bm{p}}^{(t)}\leftarrow \eta\hat{\bm{p}}^{(t)}+(1-\eta)\bm{p}^{(t)}$ \Comment{可调 EMA 平滑系数 $\eta$}
        \State $\bm{p}^{(t+1)}\leftarrow$ 对 $\tilde{\bm{p}}^{(t)}$ 施加 $p_{\min}$ 后归一化
    \EndIf
    \For{每个 $k\in \mathcal{S}_t$ \textbf{in parallel}}
        \State 接收全局模型 $w_g^{(t)}$
        \State 构造两种增强视图并计算本地对比损失
        \State 计算 joint CL/KD 与 scale CL/KD
        \State 计算周期感知尺度权重与结构约束
        \State 更新本地模型 $w_k^{(t+1)}$
    \EndFor
    \State 服务器根据样本量与质量分数聚合所有选中客户端模型
    \State 得到全局模型 $w_g^{(t+1)}$
\EndFor
\end{algorithmic}
\end{algorithm}

\section{理论分析}

本节结合本文实现，给出若干关键模块的理论解释。我们强调，这里的目标并非在当前工程实现基础上给出过强的闭式最优性证明，而是从优化与表示学习角度说明所提设计的合理性。

\subsection{非 IID 下纯参数平均的局限}

记客户端 $k$ 的局部目标为 $F_k(w)$，全局目标为
\begin{equation}
F(w)=\sum_{k=1}^{K}p_k F_k(w),\qquad p_k=\frac{n_k}{N}.
\end{equation}
若各客户端数据分布差异较大，则 $\nabla F_k(w)$ 与 $\nabla F(w)$ 之间的差异会显著增大。设存在异构上界 $\delta_k$ 使得
\begin{equation}
\|\nabla F_k(w)-\nabla F(w)\|\leq \delta_k,
\end{equation}
则局部训练 $E$ 步后，客户端相对全局模型的漂移近似满足
\begin{equation}
\|w_k^{(t+1)}-w_g^{(t)}\|
\leq \eta E\left(\|\nabla F(w_g^{(t)})\|+\delta_k\right)
 + O(\eta^2E^2).
\end{equation}
这说明当 $\delta_k$ 较大时，仅依赖局部训练和参数平均会导致显著客户端漂移\cite{fedprox,moon}。对表示学习任务而言，这种漂移不仅体现在分类精度下降上，更体现在嵌入空间结构失稳上\cite{fedu2,fedca}。

\subsection{局部--全局蒸馏对客户端漂移的抑制}

若在客户端目标中加入蒸馏约束
\begin{equation}
\mathcal{L}_{\text{KD}} = \mathbb{E}_{x\sim D_k} D\!\left(f(x;w_k), f(x;w_g)\right),
\end{equation}
其中 $D(\cdot,\cdot)$ 为非负散度，则客户端新的优化目标变为
\begin{equation}
F_k^{\text{distill}}(w)=F_k(w)+\lambda D_k(w,w_g).
\end{equation}
由于附加项在本地优化过程中显式惩罚偏离全局表示语义的解，因此新的局部解会更倾向于落在全局模型诱导的函数邻域内\cite{moon,fedx}。相比每轮直接用全局模型硬覆盖本地模型，这种做法保留了本地模型对客户端特有模式的适应能力，因此更适合时序数据的异构环境\cite{fedu2,flesd}。

\subsection{多尺度蒸馏与结构约束的作用}

设整体表示由 $R$ 个尺度组成，即
\begin{equation}
z=[z^{(1)},z^{(2)},\dots,z^{(R)}].
\end{equation}
若只对整体表示施加约束，仅能保证 $D(z_k,z_g)$ 较小，却无法排除某些尺度偏差很大、另一些尺度偏差较小而在整体上被平均掩盖的情形。因此，逐尺度约束
\begin{equation}
\sum_{r=1}^{R}\alpha_r D\!\left(z_k^{(r)}, z_g^{(r)}\right)
\end{equation}
能够提供更细粒度的误差控制。

另一方面，对尺度表示的协方差非对角项施加惩罚，本质上是在抑制不同 shapelet 通道之间的冗余，促使模型学习更加分散、更具有区分度的子模式。这有助于提高后续线性或核分类器对表示空间的利用效率。

\subsection{周期感知加权的理论动机}

设第 $r$ 个尺度的判别价值可由某种信息量指标 $I_r$ 衡量，则在固定总训练预算下，更合理的权重应满足“高价值尺度分配更大优化资源”的原则。若将总目标写为
\begin{equation}
\mathcal{L}=\sum_{r=1}^{R}\alpha_r \mathcal{L}_r,\qquad \sum_{r=1}^{R}\alpha_r=1,
\end{equation}
则当某些尺度包含更显著的周期模式时，提高其对应的 $\alpha_r$ 往往能更有效降低总体风险。第~\ref{sec:period-score} 节给出的 $\bm{\pi}$ 即此类代理的一种显式构造：STFT 分支捕获主频附近的短—中尺度能量，ACF 分支捕获长滞后峰值所暗示的长周期结构，二者经 $\beta$ 融合后得到与 shapelet 尺度对齐的权重。

\subsection{OMP 客户端选择的合理性}

本文将全局模型视为由多个候选客户端局部模型组成的“字典”所稀疏重构得到。这一建模的直觉是：在非 IID 环境下，并非所有客户端更新都同等重要，某些客户端的局部模型更接近当前全局模型所需的主要方向，因而更具有代表性\cite{secompfl,folb,updateaware}。OMP 通过求解稀疏重构问题显式寻找这种代表性组合。

若记客户端更新为 $\Delta_k$，其估计方差为 $\sigma_k^2$，则服务器聚合更新的方差近似为
\begin{equation}
\mathrm{Var}(\Delta)=\sum_{k} \omega_k^2 \mathrm{Var}(\Delta_k).
\end{equation}
若 OMP 倾向于选择对重构全局模型贡献更大的客户端，而这些客户端同时提供更稳定、更一致的更新方向，则可在有限预算下减少低质量更新进入聚合过程的机会，从而降低总体聚合方差。虽然这并不意味着 OMP 在任意场景下都严格最优，但它从“全局模型稀疏表达”的角度为客户端选择提供了一种与表示学习目标更一致的标准。指数移动平均系数 $\eta$ 与最小概率 $p_{\min}$ 分别抑制单轮系数噪声带来的采样分布震荡、并保证各客户端长期非零被采机会；二者均为实现中的\textbf{可调超参数}，可在验证集或消融实验中择优设置。

\section{实验设计}

\subsection{数据集与联邦划分}

本文实验以 UEA 多变量时间序列分类数据集为主，并进一步扩展到 HAR、Epilepsy、FD-A、SleepEDF 等数据集，以验证方法在不同规模、不同通道数和不同周期结构下的泛化能力。为模拟真实联邦环境，本文采用 Dirichlet 分布对训练集进行非 IID 划分。设客户端数量为 $K$，Dirichlet 参数为 $\alpha$，则较小的 $\alpha$ 对应更强的数据异构性。

\subsection{对比方法}

本文的客户端选择对比方法包括：
\begin{itemize}[leftmargin=2em]
    \item \textbf{uniform}：均匀随机选择客户端\cite{fedavg}；
    \item \textbf{Oort}\cite{oort}：基于客户端效用与系统速度的 guided participant selection；
    \item \textbf{FedCS}\cite{fedcsadaptive}：基于自适应采样概率的客户端选择策略；
    \item \textbf{OMP (ours)}：本文提出的基于稀疏重构的客户端选择方法，与 SecOMPFL\cite{secompfl} 同属压缩感知视角，但进一步引入 EMA 平滑与最小概率约束。
\end{itemize}

在表示学习与聚合层面，本文还通过配置开关构造若干消融版本，包括关闭 ACF 周期感知加权、关闭尺度级蒸馏、关闭联合蒸馏、关闭分布感知分数等，以分析各模块贡献。

\subsection{训练设置}

每轮训练中，服务器首先按给定策略选择客户端，再向其分发全局模型。客户端本地训练若干 epoch 后上传模型参数，服务器完成聚合。模型主干采用多尺度 shapelet 编码器，距离度量可设为混合形式；优化器采用 SGD；温度参数、损失系数、客户端数量、通信轮数、本地 epoch 数、批大小等均由配置文件统一控制。若启用 OMP 客户端选择，应在超参数表中单独列出平滑系数 $\eta$（\texttt{ema\_alpha}）与最小采样概率 $p_{\min}$（\texttt{min\_selection\_prob}）。正式实验中应给出完整超参数表，并说明不同数据集是否采用统一设置或轻度调参。

\subsection{评估协议}

本文不直接使用全局模型进行端到端分类，而是先利用全局 shapelet 编码器提取训练集和测试集表示，再在表示空间上训练 SVM 分类器。此设置与 TimeCSL 的“representation + task-oriented analyzer”范式一致\cite{timescsl}，也与 TS2Vec/TSTCC 等自监督时序表示工作的线性评测协议相容\cite{ts2vec,tstcc}，更能反映所学习表示本身的质量。评估指标包括测试准确率、平均排名、收敛轮数以及可选的训练时间指标；若后续扩展到聚类或异常检测任务，还可引入 NMI、RI 或 F1-score 等指标。

\section{结果分析与讨论}

\subsection{总体性能对比}

总体实验应重点比较以下方法：集中式 TimeCSL/CSL、uniform、Oort、FedCS、OMP，以及 OMP 与不同表示学习配置组合后的版本。预期结果是，集中式方法在数据可集中时仍然具有较强上界表现，而在联邦环境下，OMP 相较 uniform 能更有效利用有限的客户端预算；相较 Oort 和 FedCS，OMP 更强调客户端更新对全局语义结构的代表性，因此在表示学习场景下具有更好的适配性。

\begin{table*}[htbp]
  \centering
  \caption{Performance Comparison of Different Methods on Datasets }
  \label{tab:dataset_performance}
  \resizebox{\textwidth}{!}{%
  \begin{tabular}{lcccc}
    \toprule
    Dataset & Local  &  AVG &  联邦对比+多尺度自适应 \\
    \midrule
    ArticularyWordRecognition & 0.99 & 0.986666667 & 0.9833333333333333 \\
    AtrialFibrillation & 0.533 & 0.066666667 & 0.0666666666666666 \\
    BasicMotions & 1 & 0.9 & 0.95 \\
    CharacterTrajectories & 0.991 & 0.8753481894150418 & 0.867688022 \\
    Cricket & 0.994 & 0.958333333 & 0.9722222222222222 \\
    DuckDuckGeese & 0.38 & 0.24 & 0.26 \\
    EigenWorms & 0.779 & - & - \\
    Epilepsy & 0.986 & 0.934782609 & 0.949275362 \\
    ERing & 0.967 & 0.940740741 & 0.944444444444444 \\
    EthanolConcentration & 0.498 & 0.524714829 & 0.558935361 \\
    FaceDetection & 0.593 & 0.5976163450624291 & - \\
    FingerMovements & 0.59 & 0.47 & 0.53 \\
    HandMovementDirection & 0.432 & 0.469411765 & 0.486486486 \\
    Handwriting & 0.533 & 0.468235294 & 0.475294118 \\
    Heartbeat & 0.722 & 0.717073171 & 0.717073171 \\
    InsectWingbeat & 0.256 & - & - \\
    JapaneseVowels & 0.919 & 0.6081081081081081 & 0.6 \\
    Libras & 0.906 & 0.855555556 & 0.85 \\
    LSST & 0.617 & 0.5924574209245742 & 0.6042173560421735 \\
    MotorImagery & 0.61 & - & - \\
    NATOPS & 0.878 & 0.822222222 & 0.833333333 \\
    PEMS-SF & 0.827 & 0.791907514 & 0.803468208 \\
    PenDigits & 0.99 & 0.9862778730703259 & - \\
    PhonemeSpectra & 0.255 & 0.2570832090665076 & - \\
    RacketSports & 0.882 & 0.848684211 & 0.868421053 \\
    SelfRegulationSCP1 & 0.846 & 0.802047782 & 0.80887372 \\
    SelfRegulationSCP2 & 0.494 & 0.522222222 & 0.527777778 \\
    SpokenArabicDigits & 0.99 & 0.28512960436562074 & - \\
    StandWalkJump & 0.667 & 0.4 & 0.4 \\
    UWaveGestureLibrary & 0.922 & 0.921875 & 0.934375 \\
    \bottomrule
  \end{tabular}%
  }
\end{table*}

\begin{table}[htbp]
    \centering
    \caption{不同方法在各个数据集上的性能对比 (消融实验)}
    \label{tab:performance_comparison_rounded}
    \begin{tabular}{lcccc}
        \toprule
        \multirow{2}{*}{Dataset} & \multirow{2}{*}{local } & \multirow{2}{*}{ AVG} & \multirow{2}{*}{ 联邦对比All} &  联邦对比 \\
        & & & & +多尺度自适应 \\
        \midrule
        Epilepsy & & 0.9804 & 0.9839 & 0.9852 \\
        HAR      & & 0.9281 & 0.9342 & 0.9338 \\
        SleepEDF & & 0.7797 & 0.8052 & 0.8213 \\
        \bottomrule
    \end{tabular}
\end{table}

\begin{table}[htbp]
  \centering
  \caption{不同规模数据集上客户端选择的准确率表现}
  \begin{tabular}{lccccc}
    \toprule
    数据集类型 & 数据集名称 & fedcsl+omp & fedavg+uniform & Oort & FedCS \\
    \midrule
    \multirow{3}{*}{大规模数据集} & SleepEDF & 0.831987 & 0.821437 &  &  \\
    & HAR & 0.947743 & 0.927044 &  &  \\
    & HandWriting & 0.503529 & 0.490588 &  &  \\
    \midrule
    \multirow{2}{*}{中规模数据集} & PhonemeSpectra & 0.239189 & 0.211452 &  &  \\
    & SpokenArabicDigits & 0.947704 & 0.911778 &  &  \\
    \midrule
    小规模数据集 & FaceDetection & - & - &  &  \\
    & InsectWingbeat & - & - &  &  \\
    \bottomrule
  \end{tabular}
  \label{tab:scale_performance}
\end{table}
\subsection{消融实验}

消融实验用于回答“性能提升来自哪里”这一关键问题。具体而言，可依次关闭以下模块：joint CL、joint KD、scale CL、scale KD、ACF 周期感知加权和分布感知聚合。若关闭联合蒸馏后精度明显下降，则说明在非 IID 条件下特征层面对齐确实能够有效抑制客户端漂移；若关闭周期感知加权后多数据集平均精度下降，则说明不同尺度的区分性并不相同，统一等权并不是最优选择。

\subsection{非 IID 强度与收敛趋势分析}
\begin{figure}[htbp]
    \centering
    \includegraphics[width=0.95\linewidth]{../../DirichletGrid/dirichlet_alpha_grid_4x2.png}
    \caption{不同 Dirichlet 参数 $\alpha$ 下各客户端选择策略的测试准确率随 round 变化曲线（4$\times$2 子图：自上而下、自左而右依次对应从强异构到弱异构的 $\alpha$ 取值）。}
    \label{fig:dirichlet}
\end{figure}
通过改变 Dirichlet 参数 $\alpha$（对应图~\ref{fig:dirichlet}），可以分析各方法在不同异构程度下的表现。当 $\alpha$ 较小时，客户端数据分布差异更大，此时若 OMP 和局部--全局蒸馏机制依然能保持较好性能，则说明本文方法对联邦异构性具有较强鲁棒性。此外，还应给出 round 维度上的测试准确率曲线和平均损失曲线，以展示 OMP 选择与多尺度蒸馏对收敛速度和训练稳定性的影响。

\subsection{可解释性分析}

与纯 CNN 或 Transformer 表示相比，FedCSL 的一个重要优势在于 shapelet 表示的可解释性。实验中可以进一步展示：不同类别样本在表示空间中的降维分布，不同尺度 shapelet 的响应强度，以及周期感知加权下各尺度权重的变化趋势。通过这些分析，可以更直观地说明本文方法不仅提升了精度，也保留了模型的可解释性和可分析性。

\subsection{局限性讨论}

尽管本文提出了 OMP 客户端选择与联邦 shapelet 表征学习的联合框架，但仍存在若干局限。第一，OMP 优势的理论分析目前主要建立在稀疏重构和代表性更新的解释上，尚未形成严格的闭式最优性结论\cite{secompfl}。第二，本文当前实验重点仍在分类任务，聚类与异常检测等下游任务还需进一步扩展\cite{ts2vec,timescsl}。第三，系统异构因素在本文中尚未像 Oort 那样被显式建模到完整的时间代价函数中\cite{oort,fedcsadaptive}，未来可进一步融合时延感知或资源感知机制\cite{updateaware,igpe}。

\section{结论}

本文围绕联邦非 IID 多变量时间序列场景中的“表示学习”与“客户端选择”两个核心问题，提出了 FedCSL 框架。该方法以 TimeCSL/CSL 的多尺度 shapelet 表征学习为骨干，在联邦环境下进一步引入基于 OMP 的客户端选择、局部--全局联合对比与知识蒸馏、多尺度结构对齐以及周期感知加权机制。理论分析表明，所提方法能够从特征层面抑制客户端漂移、降低尺度失配并提升聚合质量；实验设计则从总体性能、消融实验、非 IID 鲁棒性和可解释性等多个角度验证其有效性。未来工作将继续完善 OMP 客户端选择的理论分析，并将该框架进一步扩展到联邦异常检测、联邦聚类和个性化联邦学习等更广泛场景中。


\begin{ack}
本节可在正式投稿前补充项目来源、基金支持、实验平台与协作者贡献等信息。匿名投稿时请勿写入可识别信息；终稿需按 NeurIPS 要求披露资助与利益冲突。
\end{ack}

% References (natbib is loaded by neurips_2023)
{\small
\bibliographystyle{unsrtnat}
\bibliography{refs}
}

\end{document}
